\documentclass[11pt]{article}
\usepackage{amsmath,amssymb,amsthm,hyperref}
\newtheorem{theorem}{Theorem}
\title{A Strict Axis-Zero Compiler from Reciprocal Sheets and Path Orders}
\author{}
\date{November 23, 2026}
\begin{document}
\maketitle

\paragraph{Status and limitation.}
This is a complete written universal local construction with a linear
head count and an improved leading constant. It does not yield an
unbounded row/head ratio, and therefore does not prove the requested
$\log\log n$ or $\log n$ saving. The original atom conditions, input
dimension, and global complexity definitions are unchanged. The proof
uses the ordinary real inverse function theorem and the Cauchy and
Borchardt determinant identities. It has not been formalized in Lean.

\begin{theorem}
For every integer $k\geq1$, every $s\geq1$, and every sign table on
$(3k-1)$ rows and $s$ columns, there is a sum of $2k+3$ affine ratios
whose denominators are positive on the entire Boolean
$(3k-1+s)$-cube and have strictly negative slopes in every coordinate.
The score has the prescribed strict signs on the one-hot product
and is exactly zero on both Boolean coordinate axes.
Consequently, for arbitrary positive $r,s$, the same conclusions hold
with at most
\[
 2\left\lceil\frac{\min(r,s)+1}{3}\right\rceil+3
\]
heads. This bound is $\frac23\min(r,s)+O(1)$, not a superconstant
compression.
\end{theorem}
\begin{proof}
Put $R=3k-1$, $m=R+s$, and $\varepsilon=1/[100(m+1)]$.
There are $2k$ regular rows indexed by $(i,t)$ with
$1\leq i\leq k$ and $t\in\{0,1\}$, and $k-1$ comparison rows
indexed by $\ell=1,\ldots,k-1$. Let
\[
 p=1-2\varepsilon,\qquad
 \lambda_j=\frac{j\varepsilon}{k+1},\qquad
 \kappa_{i,t}=\frac{i\varepsilon}{k+1},\qquad
 \kappa_\ell=\frac\varepsilon2.
\]
Thus all $\lambda_j$ and all $\kappa_\rho$ lie strictly between
zero and $\varepsilon$, and the $\lambda_j$ are distinct. Write
$d_i=p-\kappa_{i,0}$ and $d_*=p-\varepsilon/2$.

For parameters $(u,a)\in(0,\varepsilon)^k\times\mathbb R^k$
define regular queries
\[
 G_i^t(u,a)=\sum_{h=1}^k\left(
 \frac{t+a_h}{d_i-u_h}-\frac{t}{d_i}
                         -\frac{a_h}{p-u_h}\right)+c_i^t,
\]
where
\[
 c_i^0=0,\qquad
 c_i^1=-\sum_{j=1}^k\left(\frac1{d_i-\lambda_j}-\frac1{d_i}\right).
\]
Define the comparison queries
\[
 G_\ell(u,a)=\frac1{d_*-u_\ell}-\frac1{d_*-u_{\ell+1}}
       +\sum_{h=1}^k a_h\left(\frac1{d_*-u_h}-\frac1{p-u_h}\right).
\]
Every displayed denominator is positive on the parameter box.

For each permutation $\pi\in\mathfrak S_k$, let
$x_\pi$ have $u_h=\lambda_{\pi(h)}$ and $a_h=0$.
All $2k$ regular queries vanish at every $x_\pi$.
The comparison query at that point has strict sign
$\operatorname{sgn}(\lambda_{\pi(\ell)}-\lambda_{\pi(\ell+1)})$.
Every choice of the $k-1$ comparison signs is realized by a permutation:
orient the corresponding edges of the path on $k$ slots so that
an arrow points toward the slot required to have larger value.
An orientation of a path has no directed cycle. A topological ordering
therefore assigns the distinct increasing values $\lambda_j$ consistently.

It remains to check that all $2k$ regular signs can be varied
independently near each such sheet. At $x_\pi$, the Jacobian of the
regular queries in row order $(G_1^0,\ldots,G_k^0,G_1^1,\ldots,G_k^1)$
and parameter order $(u_1,\ldots,u_k,a_1,\ldots,a_k)$, after subtracting
the first block of rows from the second, is
\[
 \begin{pmatrix}0&K_\pi\\E_\pi&0\end{pmatrix},\qquad
 (K_\pi)_{ih}=\frac1{d_i-\lambda_{\pi(h)}}-
                          \frac1{p-\lambda_{\pi(h)}},\qquad
 (E_\pi)_{ih}=\frac1{(d_i-\lambda_{\pi(h)})^2}.
\]
Let $C_{ij}=1/(d_i-\lambda_j)$. The unpermuted $K$ factors as
\[
 K=\operatorname{diag}(p-d_i)\ C
             \operatorname{diag}\left(\frac1{p-\lambda_j}\right).
\]
Its two diagonal factors are positive and invertible. Cauchy's
identity gives $\det C\neq0$, and Borchardt's identity gives
$\det(C\circ C)=\det C\operatorname{perm}C\neq0$.
Hence both $K_\pi$ and $E_\pi$ are nonsingular, and so is the
full Jacobian. A primary proof of Borchardt's identity is
D.~Singer, \emph{A Bijective Proof of Borchardt's Identity},
Electronic Journal of Combinatorics 11 (2004), Theorem~1.1,
\url{https://www.combinatorics.org/ojs/index.php/eljc/article/download/v11i1r48/pdf/}.

Given all $R$ desired signs, first choose $\pi$ for its comparison
signs. Take a neighborhood of $x_\pi$ inside the parameter box on
which those strict comparison signs persist. By the inverse function
theorem, the regular queries map a sufficiently small neighborhood
of $x_\pi$ onto a neighborhood of zero in $\mathbb R^{2k}$.
This image contains a vector with any desired $2k$ strict signs.
Thus every $R$-sign vector occurs with $0<u_h<\varepsilon$.
For each column choose parameters $(u_{hj},a_{hj})$ producing that
column's desired signs.

We now realize these queries on the original Boolean cube, including
exact axis cancellation. Define affine functions
\[
 D(z,w)=1-\varepsilon\left(\sum_\rho z_\rho+\sum_jw_j\right),
 \quad U(z)=\sum_\rho\kappa_\rho z_\rho,
 \quad V_h(w)=\sum_j u_{hj}w_j,
 \quad A_h(w)=\sum_j a_{hj}w_j.
\]
For each head define $T_h(z)$ by giving it coefficient $t$ on regular
row $(i,t)$ and coefficient
$\mathbf1_{h=\ell}-\mathbf1_{h=\ell+1}$ on comparison row $\ell$.
Let $c_\rho=c_i^t$ on regular rows and $c_\rho=0$ on comparison rows.
Set
\[
 b=\frac1{p-\varepsilon}-\frac1p>0,
 \qquad L(z)=\frac1b\sum_\rho c_\rho z_\rho.
\]
The desired score is
\[
\begin{split}
 Q(z,w)={}&\sum_{h=1}^k\frac{T_h(z)+A_h(w)}{D(z,w)-U(z)-V_h(w)}
       -\frac{\sum_hT_h(z)}{D(z,w)-U(z)}\\
 &-\sum_{h=1}^k\frac{A_h(w)}{D(z,w)-V_h(w)}
       +\frac{L(z)}{D(z,w)-\varepsilon\sum_jw_j}
       -\frac{L(z)}{D(z,w)}.
\end{split}
\]
This is exactly $k+1+k+2=2k+3$ affine ratios. Each denominator
has strictly negative slopes in all coordinates, because the shared
background $D$ contributes $-\varepsilon$ everywhere and every
additional subtracted coefficient is positive. Every denominator
is at least $1-2m\varepsilon>0$ throughout the full cube: the total
subtracted coefficient sum is at most $m\varepsilon$ from $D$ and
at most another $m\varepsilon$ from its additional terms.

If $w=0$, then $V_h=A_h=0$, the first two terms cancel, and the
last two terms cancel. If $z=0$, then $U=T_h=L=0$, and the first
and third terms cancel. Thus both axis identities hold exactly.
On a one-hot grid point, $D=p$, $U=\kappa_\rho$, and the final two
terms contribute exactly $c_\rho$. Substitution gives precisely the
corresponding query $G_i^t$ or $G_\ell$ above. This proves the strict
sign claim.

For arbitrary $r$, choose $k=\lceil(r+1)/3\rceil$, pad the table
arbitrarily to $R=3k-1\geq r$ rows, and restrict the resulting
construction by fixing its extra row bits to zero. Positivity,
strict slopes on every retained coordinate, and both axis identities
are preserved. Exchanging the two blocks if necessary gives the
stated minimum of $r$ and $s$.
\end{proof}

\paragraph{Verification boundary.}
The companion exact-arithmetic checker verifies every path orientation
for $k\leq6$, the sheet values and Jacobian determinants for $k\leq4$,
and all vertices and axis identities of a small full-cube instance.
These checks validate the displayed formulas on finite examples;
the general argument is the written proof. No claim of the requested
superconstant saving or of Lean verification is made.

\paragraph{Remaining amplification problem.}
All comparison rows above use the same pole $d_*$ and at a sheet are
linear functions of the permuted vector
$(1/(d_*-\lambda_1),\ldots,1/(d_*-\lambda_k))$.
That vector lies in a fixed affine hyperplane, so this particular
family of extra readouts cannot shatter more than $k-1$ signs.
Allowing each extra row its own shared pole avoids this immediate
linear-dimensional bound, while preserving the $2k+3$ head formula.
It remains unproved that such variable-pole readouts, or more general
joint queries, can supply $\omega(k)$ independent signs.
\end{document}
